\documentclass[11pt,reqno]{amsart}

\usepackage[T1]{fontenc}
\usepackage{lmodern}
\usepackage{microtype}
\usepackage{mathtools,amssymb,amsfonts,mathrsfs}
\usepackage{booktabs}
\usepackage{xcolor}
\usepackage[margin=1.08in]{geometry}
\usepackage[colorlinks=true,linkcolor=blue!55!black,urlcolor=blue!55!black]{hyperref}

\newtheorem{theorem}{Theorem}[section]
\newtheorem{proposition}[theorem]{Proposition}
\newtheorem{lemma}[theorem]{Lemma}
\theoremstyle{remark}
\newtheorem{remark}[theorem]{Remark}

\newcommand{\R}{\mathbb R}
\newcommand{\C}{\mathbb C}
\newcommand{\one}{\mathbf 1}
\newcommand{\SO}{\mathrm{SO}}
\newcommand{\eps}{\varepsilon}
\newcommand{\ip}[2]{\left\langle #1,#2\right\rangle}
\newcommand{\abs}[1]{\left\lvert #1\right\rvert}

\title{A dimension-five counterexample to\\Makeev's conjecture}
\author{Working note}
\date{August 10, 2026}

\begin{document}

\begin{abstract}
The dimension-four counterexample suggests using the same function on
$S^4\subset\C^2\oplus\R$,
\[
 f_\eps(z_1,z_2,x)=\operatorname{Re}(z_1^2\overline z_2)
 +\eps\abs{z_1}^2\operatorname{Im}(z_1^2\overline z_2).
\]
This note proves that this function gives a counterexample in dimension five.
The last nonvanishing assertion is converted, without loss of real solutions,
to a system in a symmetric multiplication tensor.  After 14 independent
linear conditions are eliminated, the system consists of 55 quadratic
commutator equations and one cubic obstruction in 21 variables over
$\mathbb Q$.  The accompanying Magma file
\texttt{makeev\_d5\_tensor\_certificate.m} verified in 14 seconds that these
polynomials generate the unit ideal.  Thus the computer-assisted part is an
exact symbolic calculation, and the dimension-five theorem is unconditional.
\end{abstract}

\maketitle

\section{Statement}

Put
\[
 Q(z_1,z_2,x)=\operatorname{Re}(z_1^2\overline z_2),\qquad
 H(z_1,z_2,x)=\abs{z_1}^2\operatorname{Im}(z_1^2\overline z_2).
\]

\begin{theorem}\label{thm:main}
There is $\eps_0>0$ such that, whenever $0<\abs\eps<\eps_0$, the smooth odd function
$Q+\eps H$ on $S^4$ agrees with no linear functional on any rotated normalized
$A_5$ root system.  Consequently Makeev's conjecture fails in dimension five.
\end{theorem}

\section{Simplex labels and the cubic annihilator}

Let
\[
 V=\one^\perp\subset\R^6,\qquad
 q_i=e_i-\frac16\one,\qquad
 u_{ij}=\frac{q_i-q_j}{\sqrt2}.
\]
For an odd function $F$ and an isometry $g:V\to\R^5$, put
\[
 Y_F(g)_{ij}=F(gu_{ij}),\qquad Y_F(g)_{ji}=-Y_F(g)_{ij}.
\]
A skew edge-label matrix is called exact if
$Y_{ij}=\rho_i-\rho_j$.  As in the dimension-four argument,
\begin{equation}\label{eq:exact}
 Y\text{ is exact}\quad\Longleftrightarrow\quad
 \Pi Y\Pi=0,\qquad
 \Pi=I-\frac16\one\one^T.
\end{equation}
This is also equivalent to the vanishing of all triangle sums
$Y_{0i}+Y_{ij}+Y_{j0}$, $1\le i<j\le5$.

Identify $\R^5=\C^2\oplus\R$ and set
\[
 J(z_1,z_2,x)=(iz_1,2iz_2,0),\qquad
 h_\theta(z_1,z_2,x)=(e^{i\theta}z_1,e^{2i\theta}z_2,x).
\]
Both $Q$ and $H$ are invariant under $h_\theta$.  Write $p_i=gq_i$ and define
\[
 X(g)_{ij}=\ip{p_i}{Jp_j}_{\R},\qquad
 \Lambda_g(Y)=\sum_{i<j}X(g)_{ij}Y_{ij}.
\]
Since $\sum_i p_i=0$, every row sum of $X(g)$ vanishes.  Hence
$\Lambda_g$ annihilates exact labels.

The cubic identity used in dimension four is independent of the dimension.
If
\[
 P_3(v)=\sum_{i=0}^5\ip{q_i}{v}^3,
\]
then, for a homogeneous cubic $C$ and $K\in\mathfrak{so}(5)$,
\begin{equation}\label{eq:cubicidentity}
 \ip{C}{K\cdot P_3}
 =2\sqrt2\sum_{i<j}\ip{q_i}{Kq_j}_{\R}C(u_{ij}).
\end{equation}
The proof is exactly the Parseval-frame expansion from the dimension-four
paper: $\sum_iq_iq_i^T=I_V$ is the only frame fact used.  Taking
$K=g^{-1}Jg$ and $C=Q\circ g$, the infinitesimal invariance $J\cdot Q=0$
gives
\begin{equation}\label{eq:globalzero}
 \Lambda_g(Y_Q(g))=0\qquad\text{for every }g.
\end{equation}

Thus the only new statement needed in dimension five is
\begin{equation}\label{eq:target}
 Y_Q(g)\text{ exact}\quad\Longrightarrow\quad
 \Lambda_g(Y_H(g))\ne0.
\end{equation}

\section{A nine-equation description of the cubic zeros}

Write
\[
 p_i=(a_i,b_i,c_i)\in\C^2\oplus\R,
 \qquad a,b\in\C^6,\quad c\in\R^6.
\]
The frame identities say that
\[
 \one,a,b,c,\overline a,\overline b
\]
are mutually orthogonal in $\C^6$, with squared norms
$6,2,2,1,2,2$, respectively.  Equivalently,
\begin{gather*}
 \sum a_i=\sum b_i=\sum c_i=0,\qquad
 \sum a_i^2=\sum b_i^2=\sum a_ib_i=\sum a_i\overline b_i=0,\\
 \sum\abs{a_i}^2=\sum\abs{b_i}^2=2,\qquad
 \sum c_i^2=1,\qquad
 \sum c_ia_i=\sum c_ib_i=0.
\end{gather*}

At the edge $gu_{ij}$, the complex cubic has, apart from the common factor
$1/(2\sqrt2)$, the label
\[
 \widetilde Y_{ij}=(a_i-a_j)^2(\overline b_i-\overline b_j).
\]
For $(u\wedge v)_{ij}=u_iv_j-v_iu_j$, expansion modulo exact matrices gives
\begin{equation}\label{eq:omega}
 \omega:=\Pi\widetilde Y\Pi
 =-a^2\wedge\overline b+2a\wedge(a\overline b).
\end{equation}
The $Q$-labels are exact precisely when
$\operatorname{Re}\omega=0$, or $\omega+\overline\omega=0$.

Introduce
\begin{align*}
 s&=\frac12\sum\abs{a_i}^2a_i,&
 m&=\sum a_i^2\overline b_i,&
 r&=\sum c_i a_i^2,\\
 n&=\sum a_i^3,&
 \ell&=\frac12\sum a_i^2b_i,&
 \lambda&=\frac12\sum\abs{a_i}^2\overline b_i,\\
 u&=\frac12\sum a_i\overline b_i^2,&
 w&=\sum c_i a_i\overline b_i,&
 d&=\frac12\sum a_i\abs{b_i}^2.
\end{align*}
Expansion in the displayed orthogonal basis gives
\begin{align*}
 a^2&=sa+\frac m2b+rc+\frac n2\overline a+\ell\overline b,\\
 a\overline b&=\lambda a+ub+wc+\frac m2\overline a+d\overline b.
\end{align*}
Substitution in \eqref{eq:omega} yields
\begin{align*}
 \omega={}&2u\,a\wedge b+2w\,a\wedge c+m\,a\wedge\overline a
 +(2d-s)a\wedge\overline b\\
 &-\frac m2b\wedge\overline b-r\,c\wedge\overline b
 -\frac n2\overline a\wedge\overline b.
\end{align*}
The wedges of the five vectors $a,b,c,\overline a,\overline b$ are linearly
independent.  Comparing the last display with its conjugate proves the
following exact criterion.

\begin{proposition}[Cubic exactness criterion]\label{prop:criterion}
The $Q$-labels are exact if and only if
\begin{equation}\label{eq:nine}
 \operatorname{Im}m=0,\qquad
 4u=\overline n,\qquad
 r=0,\qquad w=0,\qquad 2d=s.
\end{equation}
These are nine real cubic equations.
\end{proposition}

This criterion is the key simplification: there is no need to use 25 entries
of an orthogonal matrix, nor to use a rational parametrization of
$\SO(5)$.

There is also a useful dimension-independent version.  For any number $N$ of
simplex vertices, exactness forces the components of $a^2$ and
$a\overline b$ perpendicular to
$\operatorname{span}_{\C}\{\one,a,b,\overline a,\overline b\}$ to vanish.
Together with the five coefficient equations in
\eqref{eq:nine}, this is equivalent to the two coordinatewise vector
identities
\begin{align}
 a^2&=sa+\frac m2b+\frac n2\overline a+\ell\overline b,
 \label{eq:universal-a2}\\
 a\overline b&=\lambda a+\frac{\overline n}{4}b
 +\frac m2\overline a+\frac s2\overline b.
 \label{eq:universal-ab}
\end{align}
Here $m$ is real.  Conversely, the two identities imply exactness by
substitution in \eqref{eq:omega}.  Equations
\eqref{eq:universal-a2}--\eqref{eq:universal-ab} therefore give a
quadratic, linearly growing system for testing every fixed $d\ge4$.  The file
\texttt{makeev\_general\_certificate.m} implements that system.

\section{The multiplication-tensor certificate}

For completeness, define the rational polynomial whose nonvanishing is needed.
For $d_{ij}=a_i-a_j$ and $e_{ij}=b_i-b_j$, put
\begin{align}
 \mathcal L={}&\sum_{i<j}
 \bigl(\operatorname{Im}(a_i\overline a_j)
      +2\operatorname{Im}(b_i\overline b_j)\bigr)
 \abs{d_{ij}}^2
 \operatorname{Im}(d_{ij}^2\overline e_{ij}).
 \label{eq:L}
\end{align}
All real and imaginary parts in \eqref{eq:L} are polynomials in the real
coordinates.  Directly from the normalization of the roots,
\begin{equation}\label{eq:Lscale}
 \Lambda_g(Y_H(g))=\frac{\mathcal L}{4\sqrt2}.
\end{equation}

\subsection{Encoding a simplex frame}

Write $a=x+iy$ and $b=p+iq$, with $x,y,p,q,c\in\R^6$, and put
\[
 v_0=6^{-1/2}\one,\qquad v_1=x,\quad v_2=y,\quad
 v_3=p,\quad v_4=q,\quad v_5=c.
\]
The frame identities above say precisely that $v_0,\ldots,v_5$ are an
orthonormal basis of $\R^6$.  For $1\le i,j,k\le5$, define the fully symmetric
tensor
\begin{equation}\label{eq:tensor}
 T_{ijk}=\sum_{\nu=0}^5v_i(\nu)v_j(\nu)v_k(\nu).
\end{equation}
Coordinatewise multiplication by $v_i$, expressed in this basis, has the real
symmetric matrix
\begin{equation}\label{eq:Mi}
 M_i=
 \begin{pmatrix}
  0&6^{-1/2}(\delta_{ij})_{j=1}^5\\
  6^{-1/2}(\delta_{ik})_{k=1}^5&(T_{ijk})_{j,k=1}^5
 \end{pmatrix}.
\end{equation}
Consequently
\begin{equation}\label{eq:trace-commute}
 \sum_{j=1}^5T_{ijj}=0,
 \qquad [M_i,M_j]=0\quad(1\le i<j\le5).
\end{equation}

These conditions also characterize the tensors arising from simplex frames.
Indeed, start with a real symmetric tensor satisfying
\eqref{eq:trace-commute}, and define $M_i$ by \eqref{eq:Mi}.  The entries of
the commutators outside the lower-right $5\times5$ block vanish automatically.
The matrices $M_i$ are therefore commuting and real symmetric, so they can be
simultaneously orthogonally diagonalized.  Equivalently, the commutative
algebra with unit $\sqrt6v_0$ and multiplication table
\[
 v_iv_j=\frac{\delta_{ij}}{\sqrt6}v_0+\sum_{k=1}^5T_{ijk}v_k
\]
is associative: commutation of all left-multiplication operators gives
$v_i(v_jv_k)=v_j(v_iv_k)=(v_iv_j)v_k$.  It follows that
\begin{equation}\label{eq:mult-identity}
 M_iM_j=\frac{\delta_{ij}}6I_6+\sum_{k=1}^5T_{ijk}M_k.
\end{equation}
Taking traces, and using $\operatorname{tr}M_k=0$, gives
\[
 \operatorname{tr}(M_iM_j)=\delta_{ij}.
\]
Thus the five joint eigenvalue vectors of the $M_i$ are an orthonormal basis
of $\one^\perp\subset\R^6$.  Moreover,
\[
 \operatorname{tr}(M_iM_jM_k)=T_{ijk},
\]
so simultaneous diagonalization recovers \eqref{eq:tensor}.  The tensor
description therefore introduces no spurious real frames and loses no real
frames.

For the equations used by the programs, relabel the tensor coordinates with
zero-based indices $0,1,2,3,4=x,y,p,q,c$.  The lower-right commutator entries
are
\begin{equation}\label{eq:commutator-polys}
 \sum_{m=0}^4
 \bigl(T_{ikm}T_{jm\ell}-T_{jkm}T_{im\ell}\bigr)
 +\frac{\delta_{ki}\delta_{j\ell}-\delta_{kj}\delta_{i\ell}}6=0.
\end{equation}
Here it suffices to take $i<j$ and $k<\ell$.  After zero equations and
duplicates are removed, \eqref{eq:commutator-polys} gives 55 distinct
nonzero quadratic polynomials.

\subsection{Exactness and the obstruction in tensor coordinates}

The five trace equations in \eqref{eq:trace-commute} are linear in $T$.
The nine real equations of Proposition~\ref{prop:criterion} also become
linear:
\begin{align}
 2T_{012}-T_{003}+T_{113}&=0,\notag\\
 2(T_{022}-T_{033}+2T_{123})-T_{000}+3T_{011}&=0,\notag\\
 2(T_{122}-T_{133}-2T_{023})+3T_{001}-T_{111}&=0,\notag\\
 T_{004}-T_{114}=T_{014}&=0,\notag\\
 T_{024}+T_{134}=T_{124}-T_{034}&=0,\notag\\
 2T_{022}+2T_{033}-T_{000}-T_{011}&=0,\notag\\
 2T_{122}+2T_{133}-T_{001}-T_{111}&=0.\label{eq:tensor-linear}
\end{align}
Together these are 14 independent rational linear equations in the 35
coordinates of a symmetric cubic tensor on $\R^5$.  Exact row reduction leaves
21 free tensor coordinates.

It remains to express the obstruction in those coordinates.  In addition to
$s,m,n,\ell,\lambda$ defined above, set
\[
 v=\sum_\nu a_\nu b_\nu^2,\qquad
 k=\sum_\nu b_\nu^3,\qquad
 \rho=\sum_\nu |b_\nu|^2b_\nu.
\]
All eight quantities are linear functions of $T$.  For example,
\begin{align*}
 s&=\tfrac12(T_{000}+T_{011})
       +\tfrac i2(T_{001}+T_{111}),\\
 m&=T_{002}-T_{112}+2T_{013},\\
 n&=(T_{000}-3T_{011})+i(3T_{001}-T_{111}),\\
 \ell&=\tfrac12(T_{002}-T_{112}-2T_{013})
       +\tfrac i2(T_{003}-T_{113}+2T_{012}),\\
 \lambda&=\tfrac12(T_{002}+T_{112})
       -\tfrac i2(T_{003}+T_{113}).
\end{align*}
On the linear exactness locus, $m$ is real.  Define
\begin{align*}
 \mathscr E={}&-\tfrac32m^3+\tfrac34\bar nmn
 -12\bar\rho\ell\lambda-16\bar\lambda\lambda m
 -2\bar n\bar s\ell-6\rho\lambda m
 -8\bar s\lambda n-8\lambda s^2-\bar v m v\\
 &+12\bar k\bar\lambda\ell+14\bar\ell\ell m
 +16\ell\lambda^2+4\bar\lambda\bar v n
 +5\bar\ell ns+6\bar\lambda\bar\rho m+6\bar v\ell s,
 \qquad E=\operatorname{Re}\mathscr E.
\end{align*}
A direct expansion of the edge expression \eqref{eq:L}, followed by the
exactness identities \eqref{eq:universal-a2}--\eqref{eq:universal-ab}, gives
\begin{equation}\label{eq:Escale}
 E=8\sqrt2\,\Lambda_g(Y_H(g)).
\end{equation}
Since each moment in $\mathscr E$ is linear in $T$, the polynomial $E$ is
cubic in the tensor coordinates.

\subsection{The successful exact computation}

The script \texttt{tensor\_polynomials.py} performs rational row reduction on
the 14 linear equations, substitutes the resulting parametrization into the
commutator equations and $E$, removes duplicate commutators, and writes
\texttt{makeev\_d5\_tensor\_certificate.m}.  The latter defines an ideal in
\[
 \mathbb Q[t_1,\ldots,t_{21}]
\]
with grevlex order, generated by the 55 quadratics and the cubic obstruction.
The generator clears harmless rational factors; in particular its final
polynomial is $\tfrac94E$, which has the same zero set as $E$.  It also checks
internally that exactly 56 generators were produced.

Magma was run with modular Gr\"obner-basis arithmetic enabled.  The call
\texttt{IsProper(I)} returned false, the assertion
\texttt{assert not proper} passed, and after approximately 14 seconds the
program printed
\begin{verbatim}
Certificate verified: tensor ideal is the unit ideal.
\end{verbatim}
The two source files used for this run have SHA-256 fingerprints
\begin{center}
\begin{tabular}{@{}l@{}}
\texttt{403d831600c2ac4a6d60ba768506ff031eb1d4817eadde04d80e40e3aa14d2e6}\\
\quad\texttt{makeev\_d5\_tensor\_certificate.m}\\[2pt]
\texttt{2088bc34cd12a167c9a60f83a35d24b0a61d7ef6fe61213a93a12d51d6051bc7}\\
\quad\texttt{tensor\_polynomials.py}.
\end{tabular}
\end{center}
Because every coefficient is rational, this proves that the complex affine
zero set is empty; a fortiori there is no real tensor belonging to a simplex
frame for which both \eqref{eq:nine} and $E=0$ hold.  By
\eqref{eq:Escale}, this is exactly implication \eqref{eq:target}.  Notice that
this calculation is stronger than a real nonvanishing certificate and uses no
floating-point or interval decision.

\begin{remark}
An earlier coordinate calculation used a gauge chart, 25 variables, and a
degree-seven equation after saturation.  It did not terminate in reasonable
time and is not used in the proof.  The multiplication-tensor calculation
replaces it by the successful lower-degree certificate above.
\end{remark}

\section{Completion of the proof}

If arbitrarily small nonzero $\eps$ admitted
exact labels for $Q+\eps H$, take $\eps_\nu\to0$ and corresponding rotations
$g_\nu$.  Compactness gives a subsequence $g_\nu\to g$.  Passing to the limit
in the exact-label defect shows that $Y_Q(g)$ is exact.  Since $\Lambda_{g_\nu}$
annihilates exact labels and \eqref{eq:globalzero} holds identically,
\[
 \Lambda_{g_\nu}(Y_H(g_\nu))=0.
\]
The limit contradicts \eqref{eq:target}.  This proves
Theorem~\ref{thm:main}.

An orientation issue does not arise: a transposition of two simplex vertices
induces an orientation-reversing symmetry of the $A_5$ root system.  Therefore
the $\mathrm O(5)$ and $\SO(5)$ formulations are equivalent.

\section{Evidence in higher dimensions}

The same formula makes sense on
$S^{d-1}\subset\C^2\oplus\R^{d-4}$ for every $d\ge4$, and the entire argument
through \eqref{eq:globalzero} remains valid.  Only the classification of the
cubic exact-label locus changes.

There is a distinguished cubic zero in every dimension.  With
$N=d+1$, $\zeta=e^{2\pi i/N}$, take
\[
 a_j=\sqrt{\frac2N}\,\zeta^j,\qquad
 b_j=\sqrt{\frac2N}\,\zeta^{2j}.
\]
The remaining $d-4$ real frame vectors can be completed orthonormally.  Since
\[
 (1-\zeta^r)^2(1-\overline{\zeta}^{\,2r})
\]
is purely imaginary, every $Q$-edge label is zero.  At this zero one obtains,
for $N\ge6$,
\[
 \Lambda(Y_H)=-\frac{3}{N^{3/2}}\ne0.
\]
This checks one orbit, not all of them.

The accompanying Python program searches numerically for other cubic zeros.
The following reconnaissance used independent random starts; failed starts are
excluded from the minimum.
\begin{center}
\begin{tabular}{c@{\qquad}c@{\qquad}c}
\toprule
$d$&converged starts&smallest sampled $\abs{\Lambda(Y_H)}$\\
\midrule
5&99/100&$0.037037037037\ (=1/27)$\\
6&10/10&$0.033452928407$\\
7&4/5&$0.012865396105$\\
8&8/10&$0.008659274357$\\
9&9/10&$0.009496590853$\\
10&6/10&$0.019556538182$\\
\bottomrule
\end{tabular}
\end{center}
No sampled zero had vanishing obstruction.  This is useful evidence that the
same perturbation may work beyond dimension five, but it is not a proof and
does not supply a uniform bound as $d$ varies.

For an exact test in another fixed dimension, change the first assignment
\texttt{d := 5;} in \texttt{makeev\_general\_certificate.m}.  The program
uses \eqref{eq:universal-a2}--\eqref{eq:universal-ab} and a cubic expression
for the obstruction.  A unit-ideal output proves nonvanishing in that selected
dimension; nontermination is merely an inconclusive computation.

\end{document}
